\section{Las deficiencias de la Next Token Prediction: comprimir no equivale a precisión de cálculo}

Los problemas anteriores ocurren en la visión; esta sección pasa a los modelos de lenguaje, porque Zhang Xiangyu (张祥雨) considera que en ellos también aparece un cuello de botella similar. Al entrenar modelos grandes, observó un fenómeno extraño: cuanto más grande es el modelo, más fuertes son la conversación general, la inteligencia emocional y la cantidad de conocimiento; pero las capacidades de razonamiento local, como las matemáticas, pueden primero subir, luego estancarse y, si se sigue aumentando la escala, incluso bajar. Un indicio es que los modelos grandes tienden más a «saltarse pasos», en lugar de calcular con cuidado como los modelos pequeños.

Su explicación apunta a la next token prediction. Este objetivo consiste, en esencia, en modelar la distribución de los datos, lo que también puede entenderse como compresión: cuanto mejor predice el modelo el siguiente token, mejor comprime el corpus. Pero comprimir la distribución del corpus humano no equivale a mantener la precisión de cálculo en tareas de matemáticas, lógica, geometría y física. El corpus humano contiene muchos pasos omitidos, intuiciones y elisiones; cuanto mejor comprime un modelo grande, más probable es que aprenda esos atajos.

\[
L_{\mathrm{NTP}}=-\sum_{t}\log p_{\theta}(x_t\mid x_{<t})
\]

Aquí, $x_t$ es el $t$-ésimo token; $x_{<t}$ es el contexto previo; $p_{\theta}$ es la probabilidad que el modelo asigna al siguiente token; cuanto menor es $L_{\mathrm{NTP}}$, mejor se ajusta el modelo a la distribución de entrenamiento. Pero un problema de matemáticas exige que los pasos y el resultado sean correctos, no que la distribución de la salida se parezca a las soluciones de internet.

\begin{figure}[H]
\centering
\includegraphics[width=0.96\textwidth]{next-token-compression-defect.png}
\caption{Las deficiencias de la Next Token Prediction: comprimir la distribución no equivale a precisión de cálculo; cuanto más grande es el modelo, más tiende a saltarse pasos. Diagrama conceptual propio, elaborado a partir de la conversación entre 00:41:11 y 00:50:48.}
\end{figure}

\begin{knowledgebox}{Lectura de la figura: el shortcut es un efecto secundario del aprendizaje de la distribución}
La figura va de Data Distribution a Compression y de ahí a Shortcut, Math y RL. El punto central es el shortcut: si la distribución de entrenamiento contiene muchos pasos omitidos, el modelo aprende una compresión que «se parece a las respuestas humanas», no un procedimiento que «calcula correctamente cada paso».
\end{knowledgebox}

\subsection{Colapso de características: el conflicto entre las características de alta probabilidad y las restricciones precisas}

Esta subsección explica el «colapso de características» mencionado en la entrevista. Los modelos generativos tienden a conservar las características más fáciles de comprimir, más frecuentes y más estables de la distribución de los datos; pero el razonamiento complejo suele depender de unos pocos pasos decisivos. En matemáticas, un paso equivocado invalida todo; en una imagen, si un dedo, la perspectiva o un contacto físico están mal, el conjunto pierde verosimilitud. La similitud de distribución no puede sustituir el cumplimiento de las restricciones.

\begin{figure}[H]
\centering
\includegraphics[width=0.94\textwidth]{feature-collapse.png}
\caption{El fenómeno del colapso de características: los modelos generativos conservan las características de alta probabilidad, pero pierden la precisión de cálculo y de las restricciones. Diagrama conceptual propio, elaborado a partir de la conversación entre 00:45:42 y 00:50:48.}
\end{figure}

\begin{knowledgebox}{Lectura de la figura: parecerse al corpus de entrenamiento no equivale a calcular bien}
El lado izquierdo, «distribución cercana», explica por qué el modelo tiene un lenguaje fluido y un conocimiento amplio; el lado derecho, «precisión insuficiente», explica por qué todavía se equivoca en matemáticas, geometría y física. El razonamiento multimodal necesita restricciones precisas, no solo que la salida se parezca más a las muestras de entrenamiento.
\end{knowledgebox}

\begin{warningbox}{No hay que tomar la tasa de compresión como el límite superior de la inteligencia}
Una mejor next token loss suele mejorar las capacidades generales, pero no es un objetivo suficiente para todas las tareas inteligentes. Las tareas que requieren búsqueda precisa, planificación a largo plazo, verificación externa y retroalimentación de las acciones deben incorporar objetivos y procesos adicionales.
\end{warningbox}

\subsection{Resumen de la sección}

La next token prediction es la base del despegue de los modelos grandes, pero también explica por qué un modelo más grande puede saltarse más pasos. Las matemáticas, la generación visual controlable y el razonamiento multimodal requieren funciones objetivo que vayan más allá de la compresión de la distribución.
